% =====================================================================
%  Capítulo 1: Introducción y Planteamiento del Problema
%  Rúbrica: 10 %. Debe actualizarse respecto al anteproyecto.
%  Exigido (LIN §6): contexto, antecedentes generales y técnicos, hipótesis
%  y/o preguntas de investigación, objetivos, contribuciones del proyecto.
%  Origen: v3.0 (2026-11-19). Ver 03_seleccion/triaje/2026-11-19_triaje_proyecto-titulo-v3.0.md
%  Tercera persona. Una oración por línea. Unidades con siunitx.
% =====================================================================
\chapter{Introducción y Planteamiento del Problema}
\label{cap:introduccion}

\section{Descripción del problema}
\label{sec:problema}

La Región del Maule enfrenta una severa y creciente vulnerabilidad hídrica derivada de las alteraciones del ciclo hidrológico y de prolongados eventos de sequía asociados al cambio climático~\parencite{uribe2012atlas}.
Esta situación amenaza directamente la sostenibilidad y competitividad de la producción agrícola, de manera particular el cultivo del cerezo dulce (\textit{Prunus avium} L.), una especie que ha experimentado una expansión acelerada en la zona durante el último lustro y que demanda volúmenes considerables de agua para su óptimo desarrollo~\parencite{allen1998crop}.
\revisar{La rúbrica pide justificar con datos cuantitativos y literatura actualizada: agregar cifras de déficit hídrico, superficie de cerezos en el Maule y una fuente reciente sobre la megasequía (p.\ ej.\ Garreaud et al., 2017 o 2020, verificar antes de citar).}

Ante la progresiva disminución en la disponibilidad de los recursos hídricos superficiales y subterráneos, resulta imperativo explorar y diversificar las fuentes de abastecimiento mediante la reutilización de aguas residuales tratadas.
No obstante, en zonas rurales y agrícolas descentralizadas, la implementación de sistemas de tratamiento convencionales (como lodos activados) resulta habitualmente inviable debido a sus altos costos de inversión, requerimientos energéticos y complejidad operativa~\parencite{tchobanoglous2014wastewater}.

En consecuencia, el problema a resolver consiste en la necesidad de diseñar e implementar tecnologías de depuración descentralizadas, eficientes y sostenibles ---específicamente mediante Soluciones Basadas en la Naturaleza (SbN)--- que sean capaces de tratar simultáneamente efluentes de distinta naturaleza (aguas servidas domésticas, riles de lavado de agroquímicos y aguas superficiales contaminadas).
El desafío de ingeniería radica en adaptar los modelos teóricos e hidráulicos de los humedales construidos a las solicitaciones reales de un entorno agrícola, garantizando que el efluente final cumpla con la normativa chilena de calidad para el riego agrícola y asegurando la estabilidad estructural del sistema en las instalaciones de Frutícola SABA SpA.

\section{Antecedentes}
\label{sec:antecedentes}

El desarrollo del presente proyecto de título se enmarca y fundamenta en la fase inicial de ingeniería del proyecto científico financiado por la Agencia Nacional de Investigación y Desarrollo, ANID/Anillos/ATE250033, titulado \textit{\enquote{Nature-based solution strategies for developing new sustainable water sources in sweet cherry production under a scarcity context (NbS-Waterry)}}.

Mientras que el proyecto marco contempla una investigación multidisciplinaria de tres años enfocada en evaluar la fisiología vegetal del cerezo dulce (\textit{Prunus avium} L.), la inocuidad del fruto y el Análisis de Ciclo de Vida (ACV) derivado de la reutilización de efluentes tratados, el presente trabajo se delimita al diseño de la infraestructura hidráulica y estructural.

La necesidad de este desarrollo de ingeniería surge porque los modelos teóricos de depuración mediante Soluciones Basadas en la Naturaleza (SbN), documentados ampliamente en la literatura para el tratamiento de aguas~\parencite{dotro2017treatment}, requieren ser adaptados y materializados físicamente en un entorno agrícola real bajo las exigencias normativas chilenas~\parencite{inn1987nch1333}.
El propósito de este proyecto es transformar los requerimientos biológicos e hidráulicos del sistema en un expediente técnico a nivel de ejecución.
Esto contempla establecer las memorias de cálculo, las especificaciones técnicas de materiales y la planimetría normalizada que garantizarán la posterior construcción y desarrollo de la planta piloto experimental en los terrenos de Frutícola SABA SpA.

\revisar{Agregar antecedentes técnicos de Chile: cobertura de tratamiento rural (\textcite{vera2016humedales} reporta cobertura menor al 20~\% en zonas rurales) y estado del reúso para riego, con Vera-Puerto et al. (2022), \emph{Water}, 14(4), 627, \url{https://doi.org/10.3390/w14040627} (verificado; del profesor guía): reúso directo menor al 1~\% del caudal tratado.
Modelo de redacción problema → brecha → propuesta: 05\_ejemplos/fichas/2026\_ucm\_jove\_….md.}

\section{Preguntas de investigación}
\label{sec:preguntas}

\pendiente{Exigido por los Lineamientos (§6): formular la hipótesis y/o las preguntas que guían el diseño.
Por ejemplo: ¿qué área activa de HCFV en módulos IBC satisface el balance de oxígeno para los caudales y cargas de cada afluente?
¿Qué configuración estructural de la plataforma asegura la estabilidad y la horizontalidad requerida por los sifones bajo la demanda sísmica de la NCh2369?}

\section{Objetivos}
\label{sec:objetivos}

\subsection{Objetivo general}

Diseñar la infraestructura hidráulica y estructural de una estación experimental basada en Humedales Construidos de Flujo Subsuperficial Vertical (HCFV), orientada al tratamiento y posterior reutilización agrícola de tres afluentes residuales independientes en Frutícola SABA SpA, comuna de Linares.

\subsection{Objetivos específicos}

\begin{enumerate}
  \item Dimensionar el sistema de depuración y las redes de conducción a presión y gravedad, estableciendo el balance de oxígeno de los humedales, los requerimientos operativos de los equipos de impulsión y las especificaciones técnicas para las distintas fuentes de agua.
  \item Desarrollar la modelación y el cálculo estructural de la plataforma de acero y sus fundaciones bajo solicitaciones estáticas y sísmicas normativas, integrando la planimetría consolidada y los detalles constructivos de la estación.
\end{enumerate}

\revisar{Cada objetivo agrupa varias tareas. Evaluar con el profesor guía separarlos (p.\ ej.: caracterizar caudales y cargas; dimensionar los HCFV; diseñar redes e impulsión; diseñar plataforma y fundaciones; elaborar planimetría e itemizado), para que Resultados y Conclusiones respondan uno a uno.}

\section{Contribuciones del proyecto}
\label{sec:contribuciones}

\pendiente{Exigido por los Lineamientos (§6): indicar el aporte concreto del proyecto (p.\ ej.: expediente técnico a nivel de ejecución de la primera estación experimental multifuente del proyecto ANID; adaptación del método de Platzer a módulos IBC; base para los monitoreos del proyecto marco).}

\section{Alcances y limitaciones}
\label{sec:alcances}

\pendiente{Precisar qué incluye el diseño (líneas de tratamiento, redes, plataforma, fundaciones, planos, itemizado) y qué no (operación, monitoreo, evaluación agronómica y ACV, que corresponden al proyecto marco).}
